\chapter{Situación problemática}
\label{cap:situacion}

La situación problemática se concentra en la ineficiencia del proceso de
venta y reserva de paquetes turísticos en las agencias de turismo de la
ciudad de Potosí, que afecta principalmente el control de los cupos
disponibles y la confirmación de las reservas. Como se muestra en la
Figura~\ref{fig:arbol} del Anexo~A, este problema se origina en cuatro
causas directas y produce cuatro efectos directos, cada uno con una
consecuencia indirecta.

Entre las posibles causas están la falta de un criterio uniforme para
revisar la disponibilidad antes de confirmar, la ausencia de reglas claras
para modificar los estados de reserva y la demora en comunicar cambios de
capacidad o fechas. Estos factores no dependen necesariamente del uso de
registros manuales o digitales: pueden presentarse si el proceso no está
bien definido.

De confirmarse estas dificultades, podrían ocasionar reservas mayores a la
capacidad de una salida, reprogramaciones, tiempo adicional de revisión y
reclamos por parte de los turistas. Además, la información de cotizaciones,
pagos o proveedores podría quedar descoordinada respecto de las reservas
confirmadas. Por ello, consideramos necesario estudiar primero el proceso y
luego plantear una herramienta tecnológica acorde con los resultados del
diagnóstico.
